% !TEX program = xelatex
\documentclass[aspectratio=169,11pt,handout]{beamer}
\usetheme{Madrid}
\usepackage{fontspec}
\setsansfont{DejaVu Sans}
\setmainfont{DejaVu Serif}
\usepackage{polyglossia}
\setdefaultlanguage{turkish}
\usepackage{amsmath,amssymb}
\usepackage{booktabs}
\usepackage{multicol}
\usepackage{xcolor}
\usepackage{tikz}

% ---------- Renkler ----------
\definecolor{anamavi}{RGB}{0,45,98}
\definecolor{cozumyesil}{RGB}{0,110,50}
\definecolor{uyarikirmizi}{RGB}{180,20,20}
\setbeamercolor{structure}{fg=anamavi}
\setbeamercolor{frametitle}{bg=anamavi,fg=white}
\setbeamercolor{title}{bg=anamavi,fg=white}
\setbeamercolor{palette primary}{bg=anamavi,fg=white}
\setbeamercolor{palette secondary}{bg=anamavi!85!black,fg=white}
\setbeamercolor{palette tertiary}{bg=anamavi!70!black,fg=white}
\setbeamercolor{block title}{bg=anamavi,fg=white}
\setbeamercolor{block body}{bg=anamavi!7}
\setbeamercolor{block title example}{bg=cozumyesil,fg=white}
\setbeamercolor{block body example}{bg=cozumyesil!8}
\setbeamercolor{block title alerted}{bg=uyarikirmizi,fg=white}
\setbeamercolor{block body alerted}{bg=uyarikirmizi!7}
\setbeamertemplate{navigation symbols}{}
\setbeamertemplate{blocks}[rounded][shadow=false]

% ---------- Soru / Çözüm kalıbı ----------
% \soru{başlık}{soru metni}{çözüm metni}: önce soru, bir tıklamayla çözüm
\newcommand{\soru}[3]{%
\begin{frame}{#1}
  \begin{block}{Soru}\large #2\end{block}
  \pause
  \begin{exampleblock}{Çözüm}\large #3\end{exampleblock}
\end{frame}}
\newcommand{\ab}[1]{\overline{#1}}

% ---------- Sayı doğrusu ----------
% \sd{en küçük}{en büyük}{çizim komutları}
\newcommand{\sd}[3]{\begin{tikzpicture}[scale=0.7,baseline=(current bounding box.center)]
\draw[->,thin] (#1-0.7,0)--(#2+0.7,0);
\foreach \t in {#1,...,#2}{\draw (\t,0.12)--(\t,-0.12) node[below=-1pt]{\scriptsize$\t$};}
#3\end{tikzpicture}}
\newcommand{\cizgi}[2]{\draw[anamavi,line width=2.5pt] (#1,0)--(#2,0);}
\newcommand{\isin}[2]{\draw[anamavi,line width=2.5pt,-stealth] (#1,0)--(#2,0);}
\newcommand{\acik}[1]{\draw[anamavi,thick,fill=white] (#1,0) circle (4pt);}
\newcommand{\kapali}[1]{\fill[anamavi] (#1,0) circle (4pt);}

\title[MMAT105 -- Hafta 7]{MMAT105 Genel Matematik}
\subtitle{Hafta 7: Basit Eşitsizlikler ve Mutlak Değer}
\author{}
\date{2026--2027 Güz Dönemi}

\begin{document}

% =====================================================
\begin{frame}
\titlepage
\end{frame}

% (iletişim çerçevesi kaldırıldı)

\begin{frame}{Geçen Hafta Hatırlatma: İki Bilinmeyenli Denklemler}
\begin{columns}[t]
\column{0.55\textwidth}
\textbf{Çözünüz:}
\begin{enumerate}
\item \(x+y=10,\quad x-y=4\) \ \(\Rightarrow\ (x,y)=?\)
\item \(2x+3y=12,\quad x-y=1\) \ \(\Rightarrow\ (x,y)=?\)
\item \(3x+2y=16,\quad x+2y=8\) \ \(\Rightarrow\ x\cdot y=?\)
\end{enumerate}
\column{0.4\textwidth}
\pause
\textbf{Cevaplar:}
\begin{enumerate}
\item Taraf tarafa toplama: \(2x=14\); \ \((7,3)\)
\item \(x=y+1\) yerine koyma: \(5y+2=12\); \ \((3,2)\)
\item Çıkarma: \(2x=8\Rightarrow x=4,\ y=2\); \ \(x\cdot y=8\)
\end{enumerate}
\end{columns}
\end{frame}

\begin{frame}{Bugün}
\begin{columns}[t]
\column{0.32\textwidth}
\begin{block}{Ders 1}
\begin{enumerate}
\item Eşitsizlik sembolleri, aralık gösterimi
\item Eşitsizliğin özellikleri
\item \(ax+b<0\) ve çift eşitsizlikler
\end{enumerate}
\end{block}
\column{0.32\textwidth}
\begin{block}{Ders 2}
\begin{enumerate}\setcounter{enumi}{3}
\item Tam sayı değerleri sayma
\item Bir ifadenin değer aralığı
\item \(x\pm y\), \(x\cdot y\), \(x^2\) aralıkları
\end{enumerate}
\end{block}
\column{0.32\textwidth}
\begin{block}{Ders 3}
\begin{enumerate}\setcounter{enumi}{6}
\item Mutlak değer: tanım, uzaklık
\item Mutlak değerli denklemler
\item Mutlak değerli eşitsizlikler
\end{enumerate}
\end{block}
\end{columns}
\vspace{6pt}
\begin{alertblock}{Hatırlatma}
Bu hafta ara sınav öncesi son derstir. Gelecek hafta: \textbf{ARA SINAV (Hafta 1--7 konuları)}.
\end{alertblock}
\end{frame}

% =====================================================
\section{Ders 1: Eşitsizlikler ve Aralıklar}
\begin{frame}
\centering\Large\color{anamavi}\textbf{Ders 1}\\[6pt]\normalsize Eşitsizlik Sembolleri, Aralık Gösterimi, Birinci Dereceden Eşitsizlikler
\end{frame}

\subsection{Aralık gösterimi}
\begin{frame}{Eşitsizlik Sembolleri ve Aralık Gösterimi}
\begin{columns}[t]
\column{0.44\textwidth}
\begin{block}{Semboller}
\(<\) küçük, \quad \(>\) büyük\\
\(\le\) küçük veya eşit, \quad \(\ge\) büyük veya eşit\\[3pt]
\(x\) burada bir \textbf{reel sayıdır} (aksi söylenmedikçe).
\end{block}
\begin{alertblock}{Dikkat}
\(\infty\) bir sayı değildir; \(-\infty\) ve \(\infty\) yanında parantez \textbf{daima açık}tır.
\end{alertblock}
\column{0.53\textwidth}
\begin{block}{Aralıklar (\(a<b\))}
\renewcommand{\arraystretch}{1.2}
\begin{tabular}{@{}ll@{}}
\(a<x<b\) & \((a,b)\) \ açık aralık\\
\(a\le x\le b\) & \([a,b]\) \ kapalı aralık\\
\(a\le x<b\) & \([a,b)\) \ yarı açık\\
\(a<x\le b\) & \((a,b]\) \ yarı açık\\
\(x\le a\) & \((-\infty,a]\)\\
\(x>a\) & \((a,\infty)\)
\end{tabular}\\[3pt]
Sayı doğrusunda: dahil uç \(\bullet\) (dolu), dahil olmayan uç \(\circ\) (boş).
\end{block}
\end{columns}
\end{frame}

\soru{Aralık Gösterimi -- Örnek 1}
{Aşağıdakileri aralık gösterimiyle yazınız:\\
(a) \(-2<x\le 5\) \qquad (b) \(x\ge 3\) \qquad (c) \(x<-1\) \qquad (d) \(0\le x<4\)}
{(a) \((-2,5]\) \qquad (b) \([3,\infty)\) \qquad (c) \((-\infty,-1)\) \qquad (d) \([0,4)\)\\[4pt]
Tersine: \([-1,4)\) aralığı \(\ -1\le x<4\ \) demektir; \(-1\) dahil, \(4\) dahil değil.}

\soru{Sayı Doğrusunda Gösterim -- Örnek 2}
{\((-2,3]\) aralığını ve \(x<1\) eşitsizliğini sayı doğrusunda gösteriniz.}
{\((-2,3]\): \(-2\) boş yuvarlak, \(3\) dolu yuvarlak:\quad
\sd{-3}{4}{\cizgi{-2}{3}\acik{-2}\kapali{3}}\\[6pt]
\(x<1\), yani \((-\infty,1)\):\quad
\sd{-3}{4}{\isin{1}{-3.7}\acik{1}}}

\subsection{Özellikler}
\begin{frame}{Eşitsizliğin Özellikleri}
\begin{block}{Kurallar (\(a<b\) olsun)}
\begin{enumerate}
\item Her iki tarafa aynı sayı eklenir/çıkarılır: \ \(a+c<b+c\)
\item \(c>0\) ile çarpma/bölme: yön \textbf{korunur}: \ \(ac<bc\)
\item \(c<0\) ile çarpma/bölme: yön \textbf{değişir}: \ \(ac>bc\)
\item Aynı yönlü eşitsizlikler taraf tarafa \textbf{toplanır}: \ \(a<b,\ c<d\Rightarrow a+c<b+d\)
\item \(0<a<b\ \Rightarrow\ \dfrac1a>\dfrac1b\) \quad (\(a<b<0\) için de \(\dfrac1a>\dfrac1b\))
\end{enumerate}
\end{block}
\begin{alertblock}{Sık yapılan hatalar}
Negatifle çarpıp yönü çevirmeyi unutmak. \ Eşitsizlikleri taraf tarafa \textbf{çıkarmak} veya \textbf{bölmek} (yanlış!).
\end{alertblock}
\end{frame}

\soru{Özellikler -- Örnek 3}
{\(a<b\) olduğuna göre hangileri \textbf{kesinlikle} doğrudur?\\
(I) \(a-5<b-5\) \quad (II) \(-2a<-2b\) \quad (III) \(\dfrac a3<\dfrac b3\) \quad (IV) \(4-a>4-b\)}
{(I) Doğru: iki taraftan \(5\) çıkarıldı.\\
(II) Yanlış: \(-2\) ile çarpınca yön döner: \(-2a>-2b\).\\
(III) Doğru: pozitif sayıya bölme.\\
(IV) Doğru: \(-a>-b\), sonra \(4\) eklendi.}

\soru{Özellikler -- Örnek 4 (Tuzak)}
{\(5<6\) ve \(1<4\) doğrudur. Bu iki eşitsizlik taraf tarafa toplanabilir mi, çıkarılabilir mi?}
{Toplama: \(5+1<6+4\), yani \(6<10\) \ \checkmark\\[3pt]
Çıkarma: \(5-1<6-4\), yani \(4<2\) \ \(\times\) \ \textbf{yanlış!}\\[3pt]
Aynı yönlü eşitsizlikler taraf tarafa \textbf{yalnızca toplanır}. Fark için Ders 2'deki yöntem (\(-1\) ile çarpıp toplama) kullanılır.}

\soru{Özellikler -- Örnek 5}
{\(2<x<5\) ise \(\dfrac1x\) hangi aralıktadır? \ \(-4<y<-2\) ise \(\dfrac1y\)?}
{Uçlar aynı işaretli: ters alınınca yön döner.
\[2<x<5\ \Rightarrow\ \frac15<\frac1x<\frac12\]
\[-4<y<-2\ \Rightarrow\ -\frac12<\frac1y<-\frac14\]
(Kontrol: \(y=-3\Rightarrow \frac1y=-\frac13\approx-0{,}33\).)}

\subsection{Birinci dereceden eşitsizlikler}
\begin{frame}{Birinci Dereceden Eşitsizlik Çözme}
\begin{block}{Yöntem}
Denklem gibi çözülür: bilinmeyenler bir tarafa, sayılar diğer tarafa.\\
Tek fark: \textbf{negatif} bir sayıyla çarpılır/bölünürse eşitsizliğin yönü değişir.\\[3pt]
Sonuç, çözüm kümesi (ÇK) olarak aralık gösterimiyle yazılır.
\end{block}
\begin{block}{Çift (iki taraflı) eşitsizlik}
\(m<ax+b\le n\): İşlem \textbf{üç tarafa birden} uygulanır; ortada yalnız \(x\) bırakılır.
\end{block}
\begin{alertblock}{Sık yapılan hatalar}
\(-3x<12\Rightarrow x<-4\) \ \textbf{yanlış}; \ doğrusu \(x>-4\). \quad Paydaları yok ederken her terimi çarpmayı unutmak.
\end{alertblock}
\end{frame}

\soru{Eşitsizlik Çözme -- Örnek 6}
{\(3x-5<7\) eşitsizliğinin çözüm kümesi nedir?}
{\[3x<12\ \Rightarrow\ x<4\]
\(\text{ÇK}=(-\infty,4)\)\qquad \sd{0}{6}{\isin{4}{-0.7}\acik{4}}}

\soru{Eşitsizlik Çözme -- Örnek 7}
{\(-2x+3\ge 11\) eşitsizliğinin çözüm kümesi nedir?}
{\[-2x\ge 8\]
\(-2\) ile bölünce yön değişir:
\[x\le -4\qquad \text{ÇK}=(-\infty,-4]\]}

\soru{Eşitsizlik Çözme -- Örnek 8}
{\(5-3x>2x-10\) eşitsizliğini sağlayan en büyük tam sayı kaçtır?}
{\[5+10>2x+3x\ \Rightarrow\ 15>5x\ \Rightarrow\ x<3\]
\(3\) dahil değil \(\Rightarrow\) en büyük tam sayı \(2\).}

\soru{Eşitsizlik Çözme -- Örnek 9}
{\(\dfrac{x-1}{2}-\dfrac{x+2}{3}\le 1\) eşitsizliğinin çözüm kümesi nedir?}
{Paydaları eşitlemek için her terimi \(6\) ile çarpalım (\(6>0\), yön korunur):
\[3(x-1)-2(x+2)\le 6\ \Rightarrow\ 3x-3-2x-4\le 6\ \Rightarrow\ x\le 13\]
\(\text{ÇK}=(-\infty,13]\)}

\soru{Çift Eşitsizlik -- Örnek 10}
{\(-3<2x+1\le 7\) eşitsizliğinin çözüm kümesi nedir?}
{Üç taraftan \(1\) çıkar, sonra \(2\)'ye böl:
\[-4<2x\le 6\ \Rightarrow\ -2<x\le 3\]
\(\text{ÇK}=(-2,3]\)\qquad \sd{-3}{4}{\cizgi{-2}{3}\acik{-2}\kapali{3}}}

\soru{Çift Eşitsizlik -- Örnek 11}
{\(-1\le \dfrac{5-x}{2}<3\) eşitsizliğinin çözüm kümesi nedir?}
{\[-2\le 5-x<6\ \Rightarrow\ -7\le -x<1\]
\(-1\) ile çarp, \textbf{iki yön de} değişir:
\[-1<x\le 7\qquad \text{ÇK}=(-1,7]\]}

\soru{Özel Durumlar -- Örnek 12 (Tuzak)}
{(a) \(3x-2<3x+1\) \qquad (b) \(2x+5<2(x+1)\)\\ eşitsizliklerinin çözüm kümeleri nedir?}
{(a) \(x\)'ler sadeleşir: \(-2<1\) \ her zaman doğru \(\Rightarrow\) \(\text{ÇK}=\mathbb{R}\)\\[4pt]
(b) \(2x+5<2x+2\ \Rightarrow\ 5<2\) \ hiçbir zaman doğru değil \(\Rightarrow\) \(\text{ÇK}=\varnothing\)}

\begin{frame}{Ders 1 -- Sınıfta Siz Çözün}
\begin{columns}[t]
\column{0.58\textwidth}
\begin{enumerate}
\item \(4x-7\ge 9\) \ ÇK \(=?\)
\item \(2-5x<17\) \ ÇK \(=?\) \(\star\)
\item \(1<3x-2\le 10\) \ ÇK \(=?\) \(\star\)
\item \(-3\le x<2\) aralık gösterimiyle?
\item \(3<x<6\) ise \(\dfrac1x\) hangi aralıkta? \(\star\)
\end{enumerate}
\column{0.36\textwidth}
\pause
\textbf{Cevaplar:}
\begin{enumerate}
\item \([4,\infty)\)
\item \((-3,\infty)\)
\item \((1,4]\)
\item \([-3,2)\)
\item \(\left(\dfrac16,\dfrac13\right)\)
\end{enumerate}
\end{columns}
\vspace{4pt}{\footnotesize \(\star\) işaretli soruların adım adım çözümü sonraki slaytlardadır.}
\end{frame}

% ---- Ders 1 alıştırma çözümleri ----
\soru{Ders 1 -- Alıştırma Çözümü (Soru 2)}
{\(2-5x<17\) eşitsizliğinin çözüm kümesi nedir?}
{İki taraftan \(2\) çıkar: \ \(-5x<15\).\\[3pt]
\(-5\) ile bölünce yön \textbf{değişir}:
\[x>-3\qquad \text{ÇK}=(-3,\infty)\]
\sd{-5}{2}{\isin{-3}{2.7}\acik{-3}}}

\soru{Ders 1 -- Alıştırma Çözümü (Soru 3)}
{\(1<3x-2\le 10\) eşitsizliğinin çözüm kümesi nedir?}
{Üç tarafa \(2\) ekle, sonra \(3\)'e böl (\(3>0\), yön korunur):
\[3<3x\le 12\ \Rightarrow\ 1<x\le 4\]
\(\text{ÇK}=(1,4]\)\qquad \sd{0}{5}{\cizgi{1}{4}\acik{1}\kapali{4}}}

\soru{Ders 1 -- Alıştırma Çözümü (Soru 5)}
{\(3<x<6\) ise \(\dfrac1x\) hangi aralıktadır?}
{Uçlar aynı işaretli (ikisi de pozitif): ters alınınca yön döner.
\[3<x<6\ \Rightarrow\ \frac13>\frac1x>\frac16\ \Rightarrow\ \frac16<\frac1x<\frac13\]
Kontrol: \(x=4\Rightarrow \frac14\), ve \(\frac16<\frac14<\frac13\) \checkmark}

% =====================================================
\section{Ders 2: Tam Sayı Değerleri ve Değer Aralığı}
\begin{frame}
\centering\Large\color{anamavi}\textbf{Ders 2}\\[6pt]\normalsize Tam Sayı Değerleri, Bir İfadenin Değer Aralığı, İki Değişkenli Sınırlar
\end{frame}

\subsection{Tam sayı değerleri}
\begin{frame}{Tam Sayı Değerlerini Sayma}
\begin{block}{Kural (\(a<b\) tam sayılar)}
\renewcommand{\arraystretch}{1.2}
\begin{tabular}{@{}ll@{}}
\([a,b]\) aralığındaki tam sayı sayısı: & \(b-a+1\)\\
\((a,b)\) aralığındaki tam sayı sayısı: & \(b-a-1\)\\
\([a,b)\) veya \((a,b]\) aralığındaki tam sayı sayısı: & \(b-a\)
\end{tabular}\\[3pt]
Önce eşitsizliği çöz, sonra uçların dahil olup olmadığına bak.
\end{block}
\begin{alertblock}{Sık yapılan hatalar}
Uç tam sayı değilse (ör. \(x<4{,}5\)) en büyük tam sayı \(4\)'tür. \quad
\(x<3\) ise en büyük tam sayı \(3\) \textbf{değil}, \(2\)'dir.
\end{alertblock}
\end{frame}

\soru{Tam Sayı Değerleri -- Örnek 13}
{\(2x-3<11\) eşitsizliğini sağlayan en büyük tam sayı kaçtır? En küçük doğal sayı kaçtır?}
{\[2x<14\ \Rightarrow\ x<7\]
En büyük tam sayı: \(6\).\\[3pt]
\(x<7\) koşulunu \(0\) sağlar \(\Rightarrow\) en küçük doğal sayı: \(0\).}

\soru{Tam Sayı Değerleri -- Örnek 14}
{\(-5\le 3x+1<13\) eşitsizliğini sağlayan kaç tam sayı vardır? Bunların toplamı kaçtır?}
{\[-6\le 3x<12\ \Rightarrow\ -2\le x<4\]
Tam sayılar: \(-2,-1,0,1,2,3\) \(\Rightarrow\) \(6\) tane.\\[3pt]
Toplam: \((-2)+(-1)+0+1+2+3=3\)}

\soru{Tam Sayı Değerleri -- Örnek 15}
{\(-7<2x-1<8\) eşitsizliğini sağlayan tam sayıların en küçüğü, en büyüğü ve toplamı nedir?}
{\[-6<2x<9\ \Rightarrow\ -3<x<4{,}5\]
Tam sayılar: \(-2,-1,0,1,2,3,4\) \ (\(7\) tane)\\[3pt]
En küçük \(-2\), en büyük \(4\); \ toplam: \(-2-1+0+1+2+3+4=7\)}

\soru{Tam Sayı Değerleri -- Örnek 16}
{\(\dfrac x3+2>x-4\) eşitsizliğini sağlayan kaç tane pozitif tam sayı vardır?}
{\(3\) ile çarp: \ \(x+6>3x-12\ \Rightarrow\ 18>2x\ \Rightarrow\ x<9\)\\[4pt]
Pozitif tam sayılar: \(1,2,\dots,8\) \(\Rightarrow\) \(8\) tane.}

\subsection{Değer aralığı}
\begin{frame}{Bir İfadenin Değer Aralığı}
\begin{block}{Tek değişken}
\(x\)'in aralığı biliniyorsa ifade adım adım kurulur: çarp (negatifse yön değişir), sonra ekle.
\end{block}
\begin{block}{İki değişken}
\begin{itemize}
\item \(x+y\): aralıklar taraf tarafa \textbf{toplanır}.
\item \(x-y\): önce \(y\)'nin aralığı \(-1\) ile çarpılır (\(-y\)), sonra \(x\) ile toplanır.
\item \(x\cdot y\): uç değerlerin \textbf{dört çarpımı} hesaplanır; en küçüğü ve en büyüğü alınır.
\item \(x^2\): aralık \(0\)'ı içeriyorsa en küçük değer \(0\)'dır.
\end{itemize}
\end{block}
\begin{alertblock}{Sık yapılan hatalar}
\(-3<x<2\Rightarrow 9>x^2>4\) \ \textbf{yanlış}; doğrusu \(0\le x^2<9\).
\end{alertblock}
\end{frame}

\soru{Değer Aralığı -- Örnek 17}
{\(-2<a<5\) ise \(3a-1\) hangi aralıktadır? \(3a-1\) kaç farklı tam sayı değeri alabilir?}
{\[-6<3a<15\ \Rightarrow\ -7<3a-1<14\]
Tam sayı değerleri: \(-6,-5,\dots,13\) \(\Rightarrow\) \(13-(-6)+1=20\) tane.}

\soru{Değer Aralığı -- Örnek 18}
{\(-3<x\le 2\) ise \(5-2x\) hangi aralıktadır?}
{\(-2\) ile çarp (yön değişir):
\[6>-2x\ge -4\ \Rightarrow\ -4\le -2x<6\]
\(5\) ekle: \quad \(1\le 5-2x<11\), \ yani \([1,11)\).}

\soru{İki Değişken: \(x+y\), \(x-y\) -- Örnek 19}
{\(2<x<6\) ve \(-3<y<1\) ise \(x+y\) ve \(x-y\) hangi aralıktadır?}
{Toplama: \(\ -1<x+y<7\)\\[4pt]
Fark: \(-1<-y<3\) \ (\(y\)'nin aralığı \(-1\) ile çarpıldı)
\[2<x<6\ \text{ ile toplanırsa: }\ 1<x-y<9\]}

\soru{İki Değişken: Tam Sayılar -- Örnek 20}
{\(2<x<6\), \ \(-3<y<1\) ve \(x,y\) tam sayılardır. \(x-y\)'nin en büyük değeri ve \(x+y\)'nin en küçük değeri kaçtır?}
{\(x\in\{3,4,5\}\), \quad \(y\in\{-2,-1,0\}\)\\[4pt]
\(x-y\) en büyük: \(x=5,\ y=-2\ \Rightarrow\ 7\)\\
\(x+y\) en küçük: \(x=3,\ y=-2\ \Rightarrow\ 1\)}

\soru{İki Değişken: \(x\cdot y\) -- Örnek 21}
{\(2<x<6\) ve \(-3<y<1\) ise \(x\cdot y\) hangi aralıktadır?}
{Uç değerlerin çarpımları:
\[2\cdot(-3)=-6,\quad 2\cdot 1=2,\quad 6\cdot(-3)=-18,\quad 6\cdot 1=6\]
En küçük \(-18\), en büyük \(6\) (uçlar dahil değil): \quad \(-18<x\cdot y<6\)}

\soru{Kare Aralığı -- Örnek 22}
{(a) \(-3<x<2\) \qquad (b) \(-1<x<4\)\\ ise \(x^2\) hangi aralıktadır?}
{Aralık \(0\)'ı içeriyor \(\Rightarrow\) \(x=0\) için \(x^2=0\) en küçük değer.\\[3pt]
(a) En büyük: \((-3)^2=9\) (dahil değil) \(\Rightarrow\ 0\le x^2<9\)\\
(b) En büyük: \(4^2=16\) (dahil değil) \(\Rightarrow\ 0\le x^2<16\)}

\soru{Karma -- Örnek 23}
{\(-4<x<3\), \ \(-2<y<5\) ve \(x,y\) tam sayılardır. \(x^2-y\) ifadesinin en büyük ve en küçük değeri kaçtır?}
{\(x\in\{-3,\dots,2\}\Rightarrow x^2\in\{0,1,4,9\}\); \quad \(y\in\{-1,\dots,4\}\)\\[4pt]
En büyük: \(x^2=9,\ y=-1\ \Rightarrow\ 9-(-1)=10\)\\
En küçük: \(x^2=0,\ y=4\ \Rightarrow\ 0-4=-4\)}

\begin{frame}{Ders 2 -- Sınıfta Siz Çözün}
\begin{columns}[t]
\column{0.6\textwidth}
\begin{enumerate}
\item \(-4\le 2x+2<9\) kaç tam sayı sağlar? \(\star\)
\item \(-1<a<4\) ise \(2-3a\) hangi aralıkta?
\item \(1<x<5\), \(-2<y<3\) ise \(x-y\) aralığı? \(\star\)
\item \(-5<x<2\) ise \(x^2\) aralığı?
\item \(-2<x<4\), \(1<y<6\), \(x,y\in\mathbb{Z}\): \ \(x\cdot y\) en küçük? \(\star\)
\end{enumerate}
\column{0.34\textwidth}
\pause
\textbf{Cevaplar:}
\begin{enumerate}
\item \(7\) \ (\(-3\le x<3{,}5\))
\item \((-10,5)\)
\item \((-2,7)\)
\item \([0,25)\)
\item \(-5\) \ (\(x=-1,\ y=5\))
\end{enumerate}
\end{columns}
\vspace{4pt}{\footnotesize \(\star\) işaretli soruların adım adım çözümü sonraki slaytlardadır.}
\end{frame}

% ---- Ders 2 alıştırma çözümleri ----
\soru{Ders 2 -- Alıştırma Çözümü (Soru 1)}
{\(-4\le 2x+2<9\) eşitsizliğini kaç tam sayı sağlar?}
{Üç taraftan \(2\) çıkar, sonra \(2\)'ye böl:
\[-6\le 2x<7\ \Rightarrow\ -3\le x<3{,}5\]
Tam sayılar: \(-3,-2,-1,0,1,2,3\) \(\Rightarrow\) \(7\) tane.}

\soru{Ders 2 -- Alıştırma Çözümü (Soru 3)}
{\(1<x<5\) ve \(-2<y<3\) ise \(x-y\) hangi aralıktadır?}
{Önce \(-y\): \(y\)'nin aralığını \(-1\) ile çarp (yön değişir):
\[-3<-y<2\]
\(1<x<5\) ile taraf tarafa topla:
\[1+(-3)<x-y<5+2\ \Rightarrow\ -2<x-y<7\]}

\soru{Ders 2 -- Alıştırma Çözümü (Soru 5)}
{\(-2<x<4\), \ \(1<y<6\) ve \(x,y\) tam sayılardır. \(x\cdot y\)'nin en küçük değeri kaçtır?}
{\(x\in\{-1,0,1,2,3\}\), \quad \(y\in\{2,3,4,5\}\)\\[3pt]
Çarpımın negatif olması için \(x\) negatif olmalı: tek seçenek \(x=-1\).\\[3pt]
En küçük çarpım için \(y\) en büyük: \(y=5\) \(\Rightarrow\) \(x\cdot y=-5\).}

% =====================================================
\section{Ders 3: Mutlak Değer}
\begin{frame}
\centering\Large\color{anamavi}\textbf{Ders 3}\\[6pt]\normalsize Mutlak Değer, Mutlak Değerli Denklemler ve Eşitsizlikler
\end{frame}

\subsection{Tanım ve özellikler}
\begin{frame}{Mutlak Değer: Tanım ve Özellikler}
\begin{columns}[t]
\column{0.48\textwidth}
\begin{block}{Tanım}
\[|x|=\begin{cases}x, & x\ge 0\\ -x, & x<0\end{cases}\]
\(|x|\): \(x\)'in \(0\)'a uzaklığı.\\
\(|x-a|\): \(x\) ile \(a\) arasındaki uzaklık.
\end{block}
\column{0.48\textwidth}
\begin{block}{Özellikler}
\(|x|\ge 0\), \quad \(|-x|=|x|\), \quad \(|x-y|=|y-x|\)\\[2pt]
\(|x\cdot y|=|x|\cdot|y|\), \quad \(\left|\dfrac xy\right|=\dfrac{|x|}{|y|}\ (y\neq0)\)\\[2pt]
Üçgen eşitsizliği: \(|x+y|\le|x|+|y|\).
\end{block}
\end{columns}
\begin{alertblock}{Sık yapılan hatalar}
\(|-x|\) her zaman \(x\) değildir: \(x<0\) ise \(|-x|=-x\). \quad Genelde \(|x+y|\neq|x|+|y|\): \(|3+(-5)|=2\), ama \(|3|+|-5|=8\).
\end{alertblock}
\end{frame}

\soru{Mutlak Değer -- Örnek 24}
{\(|-7|+|3|-|-2-4|\) \ ve \ \(|5-8|-|8-5|\) \ işlemlerinin sonuçları nedir?}
{\[|-7|+|3|-|-6|=7+3-6=4\]
\[|5-8|-|8-5|=3-3=0\qquad(|x-y|=|y-x|)\]}

\soru{Uzaklık Yorumu -- Örnek 25}
{Sayı doğrusunda \(3\)'e uzaklığı \(4\) birim olan sayılar hangileridir? Bunu mutlak değerle yazınız.}
{\(|x-3|=4\):
\sd{-2}{8}{\draw[<->,cozumyesil,thick] (-1,0.45)--(3,0.45);\draw[<->,cozumyesil,thick] (3,0.45)--(7,0.45);\kapali{-1}\kapali{7}\fill[uyarikirmizi] (3,0) circle (3pt);}\\[3pt]
\(x=3-4=-1\) \ veya \ \(x=3+4=7\).\\
Not: \(|x+2|=|x-(-2)|\): \(x\)'in \(-2\)'ye uzaklığıdır.}

\subsection{İşareti bilinen ifadeler}
\begin{frame}{İşareti Bilinen İfadeyi Mutlak Değerden Çıkarma}
\begin{block}{Kural}
\begin{itemize}
\item İçerisi \(\ge 0\) ise: \ \(|A|=A\) \quad (olduğu gibi çıkar)
\item İçerisi \(<0\) ise: \ \(|A|=-A\) \quad (işaretini değiştirerek çıkar)
\end{itemize}
Önce her mutlak değerin \textbf{içinin işaretine} karar ver, sonra sadeleştir.
\end{block}
\begin{alertblock}{Sık yapılan hatalar}
\(-A\) yazarken parantezi unutmak: \ \(|3-x|\) içerisi negatifse \(-(3-x)=x-3\) olur.
\end{alertblock}
\end{frame}

\soru{İşaret -- Örnek 26}
{\(a<0<b\) olduğuna göre \(|a-b|+|b|-|a|\) ifadesinin eşiti nedir?}
{\(a-b<0\Rightarrow |a-b|=b-a\); \quad \(|b|=b\); \quad \(|a|=-a\)
\[(b-a)+b-(-a)=2b\]
Kontrol: \(a=-1,\ b=2\): \ \(3+2-1=4=2\cdot 2\) \ \checkmark}

\soru{İşaret -- Örnek 27}
{\(x<-2\) olduğuna göre \(|x+2|-|3-x|\) ifadesinin eşiti nedir?}
{\(x+2<0\Rightarrow |x+2|=-x-2\); \quad \(3-x>0\Rightarrow |3-x|=3-x\)
\[(-x-2)-(3-x)=-5\]
Kontrol: \(x=-5\): \ \(3-8=-5\) \ \checkmark}

\soru{İşaret -- Örnek 28}
{\(1<x<4\) olduğuna göre \(|x-1|+|x-4|\) ifadesinin eşiti nedir?}
{\(x-1>0\), \quad \(x-4<0\):
\[(x-1)+(4-x)=3\]
Uzaklık yorumu: \(x\), \(1\) ile \(4\) arasındaysa uzaklıkları toplamı \(4-1=3\).}

\subsection{Mutlak değerli denklemler}
\begin{frame}{Mutlak Değerli Denklemler}
\begin{block}{\(|ax+b|=c\)}
\begin{itemize}
\item \(c>0\): \ \(ax+b=c\) \ veya \ \(ax+b=-c\) \quad (iki kök)
\item \(c=0\): \ \(ax+b=0\) \quad (tek kök)
\item \(c<0\): \ \(\text{ÇK}=\varnothing\) \quad (mutlak değer negatif olamaz)
\end{itemize}
\end{block}
\begin{block}{\(|f(x)|=|g(x)|\)}
\(f(x)=g(x)\) \ veya \ \(f(x)=-g(x)\)
\end{block}
\end{frame}

\soru{Denklem -- Örnek 29}
{\(|2x-3|=7\) denkleminin çözüm kümesi nedir? Kökler toplamı kaçtır?}
{\[2x-3=7\Rightarrow x=5\qquad\text{veya}\qquad 2x-3=-7\Rightarrow x=-2\]
\(\text{ÇK}=\{-2,5\}\); \quad kökler toplamı \(3\).}

\soru{Denklem -- Örnek 30 (Tuzak)}
{(a) \(|x+1|=-3\) \qquad (b) \(|3x-6|=0\) \qquad (c) \(|x-2|+5=3\)}
{(a) Mutlak değer negatif olamaz: \(\text{ÇK}=\varnothing\)\\[3pt]
(b) \(3x-6=0\Rightarrow x=2\): \(\text{ÇK}=\{2\}\)\\[3pt]
(c) \(|x-2|=-2\) \(\Rightarrow\) \(\text{ÇK}=\varnothing\)}

\soru{Denklem -- Örnek 31}
{\(|x-1|=|2x+4|\) denkleminin çözüm kümesi nedir?}
{\(x-1=2x+4\Rightarrow x=-5\)\\[2pt]
\(x-1=-(2x+4)\Rightarrow 3x=-3\Rightarrow x=-1\)\\[3pt]
\(\text{ÇK}=\{-5,-1\}\). \ Kontrol: \(|-6|=|-6|\), \ \(|-2|=|2|\) \ \checkmark}

\subsection{Mutlak değerli eşitsizlikler}
\begin{frame}{Mutlak Değerli Eşitsizlikler (\(b>0\))}
\begin{block}{Kurallar}
\begin{itemize}
\item \(|x-a|<b\ \Rightarrow\ a-b<x<a+b\) \quad (\(a\)'ya uzaklığı \(b\)'den az)
\item \(|x-a|>b\ \Rightarrow\ x<a-b\ \) veya \(\ x>a+b\)
\item \(|ax+c|<b\ \Rightarrow\ -b<ax+c<b\)
\end{itemize}
\(\le\), \(\ge\) için uçlar dahil edilir.
\end{block}
\begin{alertblock}{Sık yapılan hatalar}
\(|f(x)|<\) negatif sayı \(\Rightarrow\varnothing\); \quad \(|f(x)|\ge\) negatif sayı \(\Rightarrow\mathbb{R}\).\\
\(|x|>3\) için \(-3<x<3\) yazmak \textbf{yanlış}: doğrusu \(x<-3\) veya \(x>3\).
\end{alertblock}
\end{frame}

\soru{Eşitsizlik -- Örnek 32}
{\(|x-3|<5\) eşitsizliğinin çözüm kümesi nedir? Kaç tam sayı bu eşitsizliği sağlar?}
{\[-5<x-3<5\ \Rightarrow\ -2<x<8\]
\(\text{ÇK}=(-2,8)\); \ tam sayılar \(-1,0,\dots,7\): \ \(8-(-2)-1=9\) tane.}

\soru{Eşitsizlik -- Örnek 33}
{\(|2x+1|\le 7\) eşitsizliğini sağlayan tam sayıların toplamı kaçtır?}
{\[-7\le 2x+1\le 7\ \Rightarrow\ -8\le 2x\le 6\ \Rightarrow\ -4\le x\le 3\]
Tam sayılar \(-4,\dots,3\) (\(8\) tane): \ \(-3\), \(-2\), \(-1\) ile \(1,2,3\) birbirini götürür; toplam \(-4\).}

\soru{Eşitsizlik -- Örnek 34}
{\(|x+2|>3\) eşitsizliğinin çözüm kümesi nedir?}
{\(x+2>3\Rightarrow x>1\) \quad veya \quad \(x+2<-3\Rightarrow x<-5\)\\[3pt]
\(\text{ÇK}=(-\infty,-5)\cup(1,\infty)\)\qquad
\sd{-7}{3}{\isin{-5}{-7.7}\isin{1}{3.7}\acik{-5}\acik{1}}}

\soru{Eşitsizlik -- Örnek 35 (Tuzak)}
{(a) \(|3x-1|<-4\) \qquad (b) \(|x+7|\ge 0\) \qquad (c) \(|x-2|\le 0\)}
{(a) Mutlak değer negatif sayıdan küçük olamaz: \(\text{ÇK}=\varnothing\)\\[3pt]
(b) Her \(x\) için doğru: \(\text{ÇK}=\mathbb{R}\)\\[3pt]
(c) \(|x-2|\ge0\) olduğundan yalnızca \(|x-2|=0\): \(\text{ÇK}=\{2\}\)}

\soru{İki Taraflı -- Örnek 36}
{\(2\le|x-1|<5\) eşitsizliğini sağlayan kaç tam sayı vardır?}
{\(|x-1|<5\Rightarrow -4<x<6\)\\
\(|x-1|\ge 2\Rightarrow x\le -1\) veya \(x\ge 3\)\\[3pt]
Kesişim tam sayıları: \(-3,-2,-1,\ 3,4,5\) \(\Rightarrow\) \(6\) tane.}

\begin{frame}{Ders 3 -- Sınıfta Siz Çözün}
\begin{columns}[t]
\column{0.58\textwidth}
\begin{enumerate}
\item \(|-4|-|2-9|+|0|=?\)
\item \(a<0<b\) ise \(|a|+|b-a|-|b|=?\) \(\star\)
\item \(|2x-1|=9\) \ ÇK \(=?\)
\item \(|x+1|\le 4\) kaç tam sayı sağlar? \(\star\)
\item \(|x-2|>6\) \ ÇK \(=?\) \(\star\)
\end{enumerate}
\column{0.36\textwidth}
\pause
\textbf{Cevaplar:}
\begin{enumerate}
\item \(-3\)
\item \(-2a\)
\item \(\{-4,5\}\)
\item \(9\) \ (\(-5\le x\le 3\))
\item \((-\infty,-4)\cup(8,\infty)\)
\end{enumerate}
\end{columns}
\vspace{4pt}{\footnotesize \(\star\) işaretli soruların adım adım çözümü sonraki slaytlardadır.}
\end{frame}

% ---- Ders 3 alıştırma çözümleri ----
\soru{Ders 3 -- Alıştırma Çözümü (Soru 2)}
{\(a<0<b\) olduğuna göre \(|a|+|b-a|-|b|\) ifadesinin eşiti nedir?}
{\(a<0\Rightarrow |a|=-a\); \quad \(b-a>0\Rightarrow |b-a|=b-a\); \quad \(|b|=b\)
\[(-a)+(b-a)-b=-2a\]
Kontrol: \(a=-1,\ b=2\): \ \(1+3-2=2=-2\cdot(-1)\) \ \checkmark}

\soru{Ders 3 -- Alıştırma Çözümü (Soru 4)}
{\(|x+1|\le 4\) eşitsizliğini kaç tam sayı sağlar?}
{\[-4\le x+1\le 4\ \Rightarrow\ -5\le x\le 3\]
\([-5,3]\) kapalı aralığındaki tam sayılar: \(3-(-5)+1=9\) tane.\\[3pt]
\sd{-6}{4}{\cizgi{-5}{3}\kapali{-5}\kapali{3}}}

\soru{Ders 3 -- Alıştırma Çözümü (Soru 5)}
{\(|x-2|>6\) eşitsizliğinin çözüm kümesi nedir?}
{\(x\)'in \(2\)'ye uzaklığı \(6\)'dan büyük:\\[2pt]
\(x-2>6\Rightarrow x>8\) \quad veya \quad \(x-2<-6\Rightarrow x<-4\)\\[3pt]
\(\text{ÇK}=(-\infty,-4)\cup(8,\infty)\)\\[2pt]
\sd{-6}{10}{\isin{-4}{-6.7}\isin{8}{10.7}\acik{-4}\acik{8}}}

\begin{frame}{Özet}
\begin{columns}[t]
\column{0.32\textwidth}
\begin{block}{Ders 1}
\small
\begin{itemize}
\item \(\infty\) yanında parantez açık
\item Negatifle çarp/böl: yön değişir
\item Eşitsizlikler yalnızca toplanır
\item \(0<a<b\Rightarrow \frac1a>\frac1b\)
\end{itemize}
\end{block}
\column{0.32\textwidth}
\begin{block}{Ders 2}
\small
\begin{itemize}
\item \((a,b)\): \(b-a-1\) tam sayı
\item \(x-y\): \(-y\) aralığı ile topla
\item \(x\cdot y\): dört uç çarpımı
\item \(0\) aralıktaysa \(\min x^2=0\) (dahil)
\end{itemize}
\end{block}
\column{0.32\textwidth}
\begin{block}{Ders 3}
\small
\begin{itemize}
\item \(|x-a|\): \(a\)'ya uzaklık
\item \(|A|=c<0\Rightarrow\varnothing\)
\item \(|x-a|<b\): \(a-b<x<a+b\)
\item \(|x-a|>b\): \(x<a-b\) veya \(x>a+b\)
\end{itemize}
\end{block}
\end{columns}
\vspace{2pt}
\begin{alertblock}{Gelecek hafta: ARA SINAV (Hafta 1--7 konuları)}
Ders notundaki \textbf{karma tekrar} sorularını çözerek hazırlanın.
\end{alertblock}
\end{frame}

\begin{frame}
\centering\Large İyi çalışmalar
\end{frame}

\end{document}
